%2026_Assingment_Analysis_IA\begin{equation*}
\frac{\phi^{n+ 1}_{i+ 1, j}- 2\phi^{n+ 1}_{i, j}+ \phi^{n+ 1}_{i- 1, j}}{(\Delta x)^2}+ \frac{\phi^{n+ 1}_{i, j+ 1}- 2\phi^{n+ 1}_{i, j}+ \phi^{n+ 1}_{i, j- 1}}{(\Delta y)^2}= -\omega^{n+ 1}_{i, j}
\end{equation*}
\documentclass[dvipdfmx]{jsreport}
\usepackage{soupreamble}

\title{解析学特論IA レポート課題}
\author{M6012 川村聡一郎}
\date{}

\begin{document}
\maketitle

\begin{tcolorbox}$\fbox{1}$
    以下を示せ。
\begin{enumerate}
    \item $I, J\in \mathcal{I}_N$ならば$I\cap J\in\mathcal{I}_N, I\cup J\in\mathcal{F}_N$
    \item $I\in \mathcal{I}_N$ならば$I^c\in \mathcal{F}_N$
    \item $E_1, E_2\in \mathcal{F}_N$ならば$E_1\cap E_2, E_1\cup E_2, E_1^c\in \mathcal{F}_N$
\end{enumerate}
\end{tcolorbox}

\begin{tcolorbox}$\fbox{2}$
    $X$が有限加法族，$f$が$X$上の有限加法的測度ならば，$E_1, E_2, \ldots, E_n\in X$に対して以下が成り立つ。
\begin{enumerate}
    \item $E_1\cap E_2= (E_1^c\cup E_2^c)^c\in X$
    \item $\bigcup_{k= 1}^n E_k\in X$
    \item $f\left(\sum_{k= 1}^n E_k\right)= \sum_{k= 1}^n f(E_k)$
    \item $f\left(\bigcup_{k= 1}^n E_k\right)\le \sum_{k= 1}^n f(E_k)$
    \item $E_1\subset E_2$ならば$f(E_1)\le f(E_2)$
\end{enumerate}
\end{tcolorbox}

\begin{souanswer}
\begin{enumerate}
    \item de・Morgenの法則より${(E_1\cap E_2)}^c= {E_1}^c\cup {E_2}^c$。
    $E_1\cap E_2\stackrel{条件(2), (3)}{=} ({E_1}^c\cup {E_2}^c)^c\in X$

    \item 帰納法で示す。$n= 2$のとき条件(3)より成立。$n$のとき$E_1\cup \cdots\cup E_n\in X$と仮定する。$n+ 1$のとき条件(3)より$E_1\cup \cdots\cup E_n\cup E_{n+1}\in X$
    
    \item 帰納法で示す。$n= 2$のとき
\begin{equation*}
    (左辺)= f(E_1+ E_2)\stackrel{性質(5)}{=} f(E_1)+ f(E_2)= (右辺)
\end{equation*}
    次に$n$のとき成り立つと仮定し、$n+ 1$のとき成り立つことを示す。
\begin{equation*}
    (左辺)= f\left(\sum_{k= 1}^{n+ 1}E_k\right)= f\left(\sum_{k= 1}^n E_k+ E_{n+ 1}\right)= f\left(\sum_{k= 1}^n E_k\right)+ f(E_{n+ 1})= \sum_{k= 1}^{n+ 1} f(E_k)= (右辺)
\end{equation*}
    \item 帰納法で示す。$n= 2$のとき
\begin{equation*}
    (左辺)= f(E_1\cup E_2)= f(E_1)+ f(E_2)= (右辺)
\end{equation*}
%　まだ書き途中

    \item $E_1\subset E_2$であるから$E_2= E_1+ (E_2\setminus E_1)$とかける。すなわち$f(E_1)\le f(E_1+ (E_2\setminus E_1))= f(E_1)+ f(E_2\setminus E_1)$。ここで$E_2\setminus E_1$は空でないから$f(E_2\setminus E_1)\ge 0$。よって$E_1\subset E_2$ならば$f(E_1)\le f(E_2)$
\end{enumerate}
\end{souanswer}

\begin{tcolorbox}$\fbox{3}$
    $x= (x_1, \ldots, x_N)\in \R^N, r> 0$に対し，$B_r(x)= \left\{(y_1, \ldots, y_N)\in \R; \sum_{k= 1}^N (y_k- x_k)^2< r^2\right\}$とおく。$O\subset\R^N$が開集合であるとは，各$x\in O$に対し$R> 0$が存在して$B_r(x)\subset O$が成り立つことである。$O$を開集合，$C_r(x)$を超立方体として以下を示せ。
\begin{enumerate}
    \item $B_r(x)\subset C_r(x)\subset B_{\sqrt{N_r}}(x)$
    \item 各$x\in O$に対し，$C_r(x)\subset O$を満たす$r> 0$が存在する。
    \item $x\in O$に対し，$r_x= \sup\{r\ge 0;\ C_r(x)\subset O\}$とおく。このとき$r_x> 0, C_{r_x}\subset O$が成り立つ。
\end{enumerate}
\end{tcolorbox}

\begin{tcolorbox}$\fbox{4}$
    $\epsilon> 0, c_k\ge 0\ (1\le k\le N- 1)$に対し，$\delta= \frac{\epsilon}{\epsilon+ a+ c_1+ c_2+ \cdots+ c_{N- 1}}$とおく。このとき
\begin{equation*}
    0< \delta< 1, \delta^N+ c_1\delta^{N- 1}+ c_2\delta^{N- 2}+ \cdots+ c_{N- 1}\delta< (1+ c_1+ c_2+ \cdots+ c_{N- 1})\delta< \epsilon
\end{equation*}
    が成り立つことを示しなさい。
\end{tcolorbox}

\begin{souanswer}
    $(第一項)< (第二項)$を示す。
\begin{equation*}
    (第一項)= \delta(\delta^{N- 1}+ c_1\delta^{N- 2}+ c_2\delta^{N- 3}+ \cdots+ c_{N- 1})
\end{equation*}
    ここで$\delta$は$0< \delta< 1$であるから、$\delta^2, \delta^3, \ldots, \delta^{N- 1}< \delta$であり、
\begin{equation*}
    \delta(\delta^{N- 1}+ c_1\delta^{N- 2}+ c_2\delta^{N- 3}+ \cdots+ c_{N- 1})< (1+ c_1+ \cdots+ c_{N- 1})
\end{equation*}
    $(第二項)< (第三項)$を示す。与えられた$\delta$の定義より$(\epsilon+ 1+ c_1+ c_2+ \cdots+ c_{N- 1})\delta< \epsilon$。ここで$\epsilon> 0$であるから
\begin{equation*}
    (1+ c_1+ \cdots+ c_{N- 1})< (1+ c_1+ c_2+ \cdots+ c_{N- 1})
\end{equation*}
\end{souanswer}

\begin{tcolorbox}$\fbox{5}$
    以下の等式を示せ。
\begin{enumerate}
    \item $[([A+ (-x)]\cap E)+ x]= A\cap [E+ x], [-([-A]\cap E)]= A\cap [-E]$
    \item $\Gamma(A)= \Gamma(A\cap [E+ x])+ \Gamma(A\cap [E+ x]^c),\ \ \Gamma(A)= \Gamma(A\cap [-E])+ \Gamma(A\cap [-E]^c)$
\end{enumerate}
\end{tcolorbox}

\begin{tcolorbox}$\fbox{6}$
\begin{enumerate}
    \item $\{A_n\}$が単調増加ならば，$\bigcap_{k= n}^\infty A_k= A_n, \bigcup_{k= n}^\infty A_k= \bigcup_{k= 1}^\infty A_k\ \ (\forall n\ge 1)$が成り立つことを示せ。
    \item $\{A_n\}$が単調減少ならば$\lim_{n\to \infty} A_n= \bigcap_{n= 1}^\infty A_n$が成り立つことを示せ。
    \item $A_{2n- 1}= [0, 1], A_{2n}= [1, 2]\subset \R$に対し，$\liminf_{n\to \infty} A_n, \limsup_{n\to \infty} A_n$を求め，これらが一致しないことを示せ。
\end{enumerate}
\end{tcolorbox}

\begin{souanswer}
\begin{enumerate}
    \item 
\begin{enumerate}
    \item $\bigcap_{k= n}^\infty A_k= A_n$を示す。
\begin{equation*}
    \bigcap_{k= n}^\infty A_k= A_n\cap A_{n+ 1}\cap \cdots= A_n\ \ (\because 単調性)
\end{equation*}
    \item $\bigcup_{k= n}^\infty A_k= \bigcup_{k= 1}^\infty A_k$を示す。
\begin{equation*}
    \bigcup_{k= 1}^\infty A_k- \bigcup_{k= n}^\infty= \bigcup_{k= 1}^{n- 1} A_k
\end{equation*}
    ここで$\{A_k\}$は単調増加であるから、$\bigcup_{k= 1}^{n- 1} A_k\subset A_n$
\end{enumerate}
    \item 
\begin{align*}
    (左辺)
    &= \liminf_{n\to \infty}A_n= \limsup_{n\to \infty}A_n\\
    &= \bigcup_{n= \infty}^\infty \bigcap_{k= n}^\infty A_n= \bigcap_{n= 1}^\infty \bigcup_{k= n}^\infty A_n\\
    &= \bigcap_{n= 1}^\infty A_n\\
    &= (右辺)
\end{align*}
\end{enumerate}
\end{souanswer}

\begin{tcolorbox}$\fbox{7}$
\begin{enumerate}
    \item $\exists n_0\in \N \st\mu\left(\bigcup_{n= n_0}^\infty E_n\right)< \infty$ならば$\mu\left(\limsup_{n\to \infty} E_n\right)\ge \limsup_{n\to \infty} \mu(E_n)$
\end{enumerate}
\end{tcolorbox}

\begin{tcolorbox}$\fbox{8}$
    $E\subset \R^N, f\colon E\to [-\infty, \infty], a\in \R$に対し，以下を示せ。
\begin{enumerate}
    \item $\{r_n\}$を$a$に収束する狭義単調現象列とするとき，$E(f> a)= \bigcup_{n= 1}^\infty E(f> r_n)$が成り立つ
    \item $\{r_n\}$を$a$に収束する狭義単調増加列とするとき，$E(f\ge a)= \bigcap_{n= 1}^\infty E(f> r_n)$が成り立つ。
    \item $E(f\ge a)= \bigcap_{n= 1}^\infty E\left(f> a- \frac{1}{n}\right)$
    \item $E(f< +\infty)= \bigcup_{n= 1}^\infty E(f< n), E(f> -\infty)= \bigcup_{n= 1}^\infty E(f> -n)$
\end{enumerate}

\end{tcolorbox}

\begin{tcolorbox}$\fbox{9}$
    $f\colon E\to [-\infty, \infty]$の値域を$\{a_1, \ldots, a_n\}$とし，$E_k= E(f= a_k)$とおく。このとき，以下を示しなさい。
\begin{enumerate}
    \item $E_k\ (1\le k\le n)$は互いに交わらない。また$E= \sum_{k= 1}^n E_k$が成り立つ。
    \item $f= \sum_{k= 1}a_k\chi_{E_k}$が成り立つ。
\end{enumerate}
\end{tcolorbox}

\begin{tcolorbox}$\fbox{10}$
    以下を示せ。
\begin{enumerate}
    \item 集合$A, B$に対し，$\chi_A(x)\chi_B(x)= \chi_{A\cap B}(x)$が成り立つ。
    \item 互いに交わらない集合$A, B$に対し，$\chi_A(x)+ \chi_B(x)=\chi_{A+ B}(x)$が成り立つ。
    \item $E$上の2つの単関数$f= \sum_{k= 0}^m \alpha_k \chi_{E_k},\ g= \sum_{l= 0}^n \beta_l\chi_{F_l}, E= \sum_{k= 0}^m E_k= \sum_{l= 0}^n F_l$に対し，$f+ g= \sum_{k= 0}^m\sum_{l= 0}^n (\alpha_k+ \beta_l)\chi_{E_k+ F_l}$である。
\end{enumerate}
\end{tcolorbox}

\begin{tcolorbox}$\fbox{11}$
    複素数値可測関数$f, g$に対して，以下が成り立つ。
\begin{enumerate}
    \item $E(f= g)\in \mathcal{M}_N$
    \item $\alpha f+ \beta g(\alpha, \beta\in \C), fg, \bar{f}, |f|$は可測である。
    \item $f_N\colon E\to \C\ (n= 1, 2, 3, \ldots)$は可測で，各$x\in E$に対し有限な極限$\lim_{n\to \infty} f_n(x)$が存在するとき，$\lim_{n\to \infty} f_n\colon E\to \C$である。
\end{enumerate}
\end{tcolorbox}

\begin{tcolorbox}$\fbox{12}$


\end{tcolorbox}

\begin{tcolorbox}$\fbox{13}$
\begin{enumerate}
    \item $\R^N$の空でない集合$A$に対して，関数$d_A\colon \R^N\to \R$を$d_A(x)= \mathrm{dist}(x, A)= \inf_{y\in A}|x- y|\ (x\in \R^N)$で定める。このとき以下を示せ。
\begin{enumerate_alpha}
    \item $d_A(x)= 0\Leftrightarrow x\in \bar{A}$
    \item $|d_A(x_1)- d_A(x_2)|\le |x_1- x_2|\ (x_1, x_2\in \R^N)$
\end{enumerate_alpha}
    \item $F, G$はそれぞれ$\R^N$の閉集合，開集合で，$F\subset G$を満たすとする。このとき以下の問に答えよ。
\begin{enumerate_alpha}
    \item 任意の$x\in \R^N$に対し，$d_F(x)+ d_{G^c}(x)> 0$が成り立つ。
    \item $g(x)= \frac{d_{G^c}(x)}{d_F(x)+ d_{G^c}(x)}$は$\R^N$上の連続関数で，$0\le g(x)\le 1, g(x)= \begin{cases}1\ \ (x\in F)\\0\ \ (x\in G^c)\end{cases}$を満たす。
\end{enumerate_alpha}
\end{enumerate}
\end{tcolorbox}

\begin{tcolorbox}$\fbox{14}$
    非負値可積分関数$f$に対して定理3.9の$f_\epsilon$が存在することを用いて，一般の可積分関数$f\colon E\to [-\infty, \infty]$ならびに$f\colon E\to \C$に対して定理3.9の$f_\epsilon$が存在することを示せ。
\end{tcolorbox}

\begin{tcolorbox}$\fbox{15}$
\begin{enumerate}
    \item $A\subset \R^N, y\in \R^N$に対し，$(\tau_y\chi_A)(x)= \chi_{[A+ y]}(x)$が成り立つことを示せ。
    \item $\R^N$上の関数$f$と$a> 0$に対し，$E_a= \{x\in \R^n; f(x)> a\}$とおく。このとき
\begin{equation*}
    \{x\in \R^N; \tau_yf(x)> a\}= [E_a+ y]\ \ (y\in \R^N),\ \ \{x\in \R^N; \bar{f}(x)> a\}= [-E_a]
\end{equation*}
    が成り立つことを示せ。
    \item 定理3.10を証明の方針に従って証明せよ。
\end{enumerate}    
\end{tcolorbox}

\begin{tcolorbox}$\fbox{16}$
\begin{enumerate}
    \item 定理4.5'の証明の方針の$f^R_n, f^I_n$に対し、$|f^R_n(x)|\le \phi(x), |f^I_n(x)|\le \phi(x)\ \ (x\in E)$が成り立つことを示せ。
    \item 定理4.5'を証明の方針に従って示せ。
\end{enumerate}
\end{tcolorbox}

\begin{tcolorbox}$\fbox{17}$
    定理4.5''を示せ。
\end{tcolorbox}

\begin{tcolorbox}$\fbox{18}$
    Fubiniの定理IIの(1-c)並びに(2),(3)の主張を、Fubiniの定理Iと補題5.12、5.13を用いて証明せよ。
\end{tcolorbox}

\begin{tcolorbox}$\fbox{19}$
    以下を示せ。
\begin{enumerate}
    \item $\{\phi_n\}$は非負値可測単関数の増加列。$\{\psi_n\}$は非負値可測単関数の減少列である。
    \item 各$x\in [a, b]$に対し、$\phi(x)= \lim_{n\to \infty} \phi_n(x), \psi(x)= \lim_{n\to \infty} \psi_n(x)$が存在する。
    \item $\phi, \psi$は可測関数である。
    \item $0\le \phi_n(x)\le \phi(x)\le f(x)\le \psi(x)\le \psi_n(x)\le M\ \ (x\in [a, b])$
    \item $\lim_{n\to \infty}s_n= \lim_{n\to \infty}S_n= \int_a^b f(x)dx$
    \item $\lim_{n\to \infty}\int_{[a, b]} \phi_n d\mu= \int_{[a, b]} \phi d\mu,\ \ \lim_{n\to \infty} \int_{[a, b]}\psi_nd\mu= \int_{[a, b]}\psi d\mu$
\end{enumerate}
\end{tcolorbox}

\begin{tcolorbox}$\fbox{20}$
    

\end{tcolorbox}

\begin{tcolorbox}$\fbox{21}$
    

\end{tcolorbox}

\begin{tcolorbox}$\fbox{22}$
    $f_n(x)= \sum_{k= 0}^n (-1)^{k}x^{2k}\ \ (n= 0, 1, 2, \ldots), f(x)= \frac{1}{1+ x^2}$とおくとき、以下を示しなさい。
\begin{enumerate}
    \item 任意の$x\in [0, 1)$に対し、$|f_n(x)|\le \frac{2}{1+ x^2}$が成り立つ。
    \item 任意の$x\in [0, 1)$に対し、$\lim_{n\to \infty} f_n(x)= f(x)$が成り立つ。
    \item $\int_{[0, 1)} fd\mu= \lim_{n\to \infty}\int_{[0, 1)} f_nd\mu$が成り立つ。
    \item $\sum_{n= 0}^\infty \frac{(-1)^n}{2n+ 1}= \frac{\pi}{4}$が成り立つ。
\end{enumerate}

\end{tcolorbox}

\begin{tcolorbox}$\fbox{23}$
    $f, g$を$\R$上の可積分関数とし、$h(x)= \int_{-\infty}^{\infty} f(x- y)g(y)dy$とおく。
\begin{enumerate}
    \item $h$が$\R$上可積分であること、並びに$\int_{-\infty}^{\infty} f(x)dx\int_{-\infty}^{\infty} g(x)dx$が成り立つことを示しなさい。
    \item $\R$上の可積分関数$\phi$に対し、$\R$上の関数$\hat{\phi}$を$\hat{\phi}(\xi)= \int_{-\infty}^{\infty} \phi(x)e^{-ix\xi} dx$で定義する。このとき、$\hat{h}(\xi)= \hat{f}(\xi)\hat{g}(\xi)$が成り立つことを示せ。
\end{enumerate}
\end{tcolorbox}

\end{document}

\begin{equation*}
%\begin{split*}
    \omega^{n+ 1}_{i, j}= \omega^n_{i, j}+ \Delta t
    \left\{
    - \left( u^n_{i, j}\frac{\omega^n_{i+ 1, j}- \omega^n_{i- 1, j}}{2\Delta x}
    - |u_{i, j}|\frac{\omega^n_{i+ 1, j}- 2\omega^n_{i, j}+ \omega^n_{i- 1, j}}{2\Delta x}\right)
    -\left(v_{i, j}\frac{\omega^n_{i, j+ 1}- \omega^n_{i, j- 1}}{2\Delta y}
    - |v_{i, j}|\frac{\omega^n_{i, j+ 1}- 2\omega^n_{i, j}+ \omega^n_{i, j- 1}}{2\Delta y}\right)
    + \frac{1}{\mathrm{Re}}\left(\frac{\omega^n_{i+1, j}- 2\omega^n_{i, j}+ \omega^n_{i- 1, j}}{(2\Delta x)^2}+ \frac{\omega^n_{i, j+ 1}- 2\omega^n_{i, j}+ \omega^n_{i, j- 1}}{(2\Delta y)^2}\right)
    \right\}
%\end{split*}
\end{equation*}

\begin{equation*}
    \frac{\phi^{n+ 1}_{i+ 1, j}- 2\phi^{n+ 1}_{i, j}+ \phi^{n+ 1}_{i- 1, j}}{(\Delta x)^2}+ \frac{\phi^{n+ 1}_{i, j+ 1}- 2\phi^{n+ 1}_{i, j}+ \phi^{n+ 1}_{i, j- 1}}{(\Delta y)^2}= -\omega^{n+ 1}_{i, j}
\end{equation*}























































\end{document}